%%%%%%%%%%%%%%%%%%%%%%%%%%%%%%%%%%%%%%%%%%%%%%%%%%%%%%%%%%%%%%%%%%%%%%%%
%%%%%%%%%%%%%%%%%%%%%% Simple LaTeX CV Template %%%%%%%%%%%%%%%%%%%%%%%%
%%%%%%%%%%%%%%%%%%%%%%%%%%%%%%%%%%%%%%%%%%%%%%%%%%%%%%%%%%%%%%%%%%%%%%%%

%%%%%%%%%%%%%%%%%%%%%%%%%%%%%%%%%%%%%%%%%%%%%%%%%%%%%%%%%%%%%%%%%%%%%%%%
%% NOTE: If you find that it says                                     %%
%%                                                                    %%
%%                           1 of ??                                  %%
%%                                                                    %%
%% at the bottom of your first page, this means that the AUX file     %%
%% was not available when you ran LaTeX on this source. Simply RERUN  %%
%% LaTeX to get the ``??'' replaced with the number of the last page  %%
%% of the document. The AUX file will be generated on the first run   %%
%% of LaTeX and used on the second run to fill in all of the          %%
%% references.                                                        %%
%%%%%%%%%%%%%%%%%%%%%%%%%%%%%%%%%%%%%%%%%%%%%%%%%%%%%%%%%%%%%%%%%%%%%%%%

%%%%%%%%%%%%%%%%%%%%%%%%%%%% Document Setup %%%%%%%%%%%%%%%%%%%%%%%%%%%%

% Don't like 10pt? Try 11pt or 12pt
\documentclass[10pt]{article}

% The automated optical recognition software used to digitize resume
% information works best with fonts that do not have serifs. This
% command uses a sans serif font throughout. Uncomment both lines (or at
% least the second) to restore a Roman font (i.e., a font with serifs).
%\usepackage{times}
%\renewcommand{\familydefault}{\sfdefault}

% This is a helpful package that puts math inside length specifications
\usepackage{calc}
\usepackage{comment}
\usepackage{graphicx}

% Simpler bibsection for CV sections
% (thanks to natbib for inspiration)
\makeatletter
\newlength{\bibhang}
\setlength{\bibhang}{1em} %1em}
\newlength{\bibsep}
 {\@listi \global\bibsep\itemsep \global\advance\bibsep by\parsep}
\newenvironment{bibsection}%
        {\begin{enumerate}{}{%
%        {\begin{list}{}{%
       \setlength{\leftmargin}{\bibhang}%
       \setlength{\itemindent}{-\leftmargin}%
       \setlength{\itemsep}{\bibsep}%
       \setlength{\parsep}{\z@}%
        \setlength{\partopsep}{0pt}%
        \setlength{\topsep}{0pt}}}
        {\end{enumerate}\vspace{-.6\baselineskip}}
%        {\end{list}\vspace{-.6\baselineskip}}
\makeatother

% Layout: Puts the section titles on left side of page
\reversemarginpar

%
%         PAPER SIZE, PAGE NUMBER, AND DOCUMENT LAYOUT NOTES:
%
% The next \usepackage line changes the layout for CV style section
% headings as marginal notes. It also sets up the paper size as either
% letter or A4. By default, letter was used. If A4 paper is desired,
% comment out the letterpaper lines and uncomment the a4paper lines.
%
% As you can see, the margin widths and section title widths can be
% easily adjusted.
%
% ALSO: Notice that the includefoot option can be commented OUT in order
% to put the PAGE NUMBER *IN* the bottom margin. This will make the
% effective text area larger.
%
% IF YOU WISH TO REMOVE THE ``of LASTPAGE'' next to each page number,
% see the note about the +LP and -LP lines below. Comment out the +LP
% and uncomment the -LP.
%
% IF YOU WISH TO REMOVE PAGE NUMBERS, be sure that the includefoot line
% is uncommented and ALSO uncomment the \pagestyle{empty} a few lines
% below.
%

%% Use these lines for letter-sized paper
\usepackage[paper=letterpaper,
            %includefoot, % Uncomment to put page number above margin
            marginparwidth=1.2in,     % Length of section titles
            marginparsep=.05in,       % Space between titles and text
            margin=1in,               % 1 inch margins
            includemp]{geometry}

%% Use these lines for A4-sized paper
%\usepackage[paper=a4paper,
%            %includefoot, % Uncomment to put page number above margin
%            marginparwidth=30.5mm,    % Length of section titles
%            marginparsep=1.5mm,       % Space between titles and text
%            margin=25mm,              % 25mm margins
%            includemp]{geometry}

%% More layout: Get rid of indenting throughout entire document
\setlength{\parindent}{0in}

\usepackage[shortlabels]{enumitem}

%% Reference the last page in the page number
%
% NOTE: comment the +LP line and uncomment the -LP line to have page
%       numbers without the ``of ##'' last page reference)
%
% NOTE: uncomment the \pagestyle{empty} line to get rid of all page
%       numbers (make sure includefoot is commented out above)
%
\usepackage{fancyhdr,lastpage}
\pagestyle{fancy}
%\pagestyle{empty}      % Uncomment this to get rid of page numbers
\fancyhf{}\renewcommand{\headrulewidth}{0pt}
\fancyfootoffset{\marginparsep+\marginparwidth}
\newlength{\footpageshift}
\setlength{\footpageshift}
          {0.5\textwidth+0.5\marginparsep+0.5\marginparwidth-2in}
\rhead{Shun Zhang, Computer Science}
\lfoot{\hspace{\footpageshift}%
       \parbox{4in}{\, \hfill %
                    \arabic{page} of \protect\pageref*{LastPage} % +LP
%                    \arabic{page}                               % -LP
                    \hfill \,}}

% Finally, give us PDF bookmarks
\usepackage{color,hyperref}
\definecolor{darkblue}{rgb}{0.0,0.0,0.3}
\hypersetup{colorlinks,breaklinks,
            linkcolor=darkblue,urlcolor=darkblue,
            anchorcolor=darkblue,citecolor=darkblue}

%%%%%%%%%%%%%%%%%%%%%%%% End Document Setup %%%%%%%%%%%%%%%%%%%%%%%%%%%%


%%%%%%%%%%%%%%%%%%%%%%%%%%% Helper Commands %%%%%%%%%%%%%%%%%%%%%%%%%%%%

% The title (name) with a horizontal rule under it
% (optional argument typesets an object right-justified across from name
%  as well)
%
% Usage: \makeheading{name}
%        OR
%        \makeheading[right_object]{name}
%
% Place at top of document. It should be the first thing.
% If ``right_object'' is provided in the square-braced optional
% argument, it will be right justified on the same line as ``name'' at
% the top of the CV. For example:
%
%       \makeheading[\emph{Curriculum vitae}]{Your Name}
%
% will put an emphasized ``Curriculum vitae'' at the top of the document
% as a title. Likewise, a picture could be included:
%
%   \makeheading[\includegraphics[height=1.5in]{my_picutre}]{Your Name}
%
% the picture will be flush right across from the name.
\newcommand{\makeheading}[2][]%
        {\hspace*{-\marginparsep minus \marginparwidth}%
         \begin{minipage}[t]{\textwidth+\marginparwidth+\marginparsep}%
             {\large \bfseries #2 \hfill #1}\\[-0.15\baselineskip]%
                 \rule{\columnwidth}{1pt}%
         \end{minipage}}

% The section headings
%
% Usage: \section{section name}
\renewcommand{\section}[1]{\pagebreak[3]%
    \hyphenpenalty=10000%
    \vspace{1.3\baselineskip}%
    \phantomsection\addcontentsline{toc}{section}{#1}%
    \noindent\llap{\scshape\smash{\parbox[t]{\marginparwidth}{\raggedright #1}}}%
    \vspace{-\baselineskip}\par}

% An itemize-style list with lots of space between items
\newenvironment{outerlist}[1][\enskip\textbullet]%
        {\begin{itemize}[#1,leftmargin=*]}{\end{itemize}%
         \vspace{-.6\baselineskip}}

% An environment IDENTICAL to outerlist that has better pre-list spacing
% when used as the first thing in a \section
\newenvironment{lonelist}[1][\enskip\textbullet]%
        {\begin{list}{#1}{%
        \setlength{\partopsep}{0pt}%
        \setlength{\topsep}{0pt}}}
        {\end{list}\vspace{-.6\baselineskip}}

% An itemize-style list with little space between items
\newenvironment{innerlist}[1][\enskip\textbullet]%
        {\begin{itemize}[#1,leftmargin=*,parsep=0pt,itemsep=0pt,topsep=0pt,partopsep=0pt]}
        {\end{itemize}}

% An environment IDENTICAL to innerlist that has better pre-list spacing
% when used as the first thing in a \section
\newenvironment{loneinnerlist}[1][\enskip\textbullet]%
        {\begin{itemize}[#1,leftmargin=*,parsep=0pt,itemsep=0pt,topsep=0pt,partopsep=0pt]}
        {\end{itemize}\vspace{-.6\baselineskip}}

% To add some paragraph space between lines.
% This also tells LaTeX to preferably break a page on one of these gaps
% if there is a needed pagebreak nearby.
\newcommand{\blankline}{\quad\pagebreak[3]}
\newcommand{\halfblankline}{\quad\vspace{-0.5\baselineskip}\pagebreak[3]}

% Uses hyperref to link DOI
\newcommand\doilink[1]{\href{http://dx.doi.org/#1}{#1}}
\newcommand\doi[1]{doi:\doilink{#1}}

% For \url{SOME_URL}, links SOME_URL to the url SOME_URL
\providecommand*\url[1]{\href{#1}{#1}}
% Same as above, but pretty-prints SOME_URL in teletype fixed-width font
\renewcommand*\url[1]{\href{#1}{\texttt{#1}}}

% For \email{ADDRESS}, links ADDRESS to the url mailto:ADDRESS
\providecommand*\email[1]{\href{mailto:#1}{#1}}
% Same as above, but pretty-prints ADDRESS in teletype fixed-width font
%\renewcommand*\email[1]{\href{mailto:#1}{\texttt{#1}}}

%\providecommand\BibTeX{{\rm B\kern-.05em{\sc i\kern-.025em b}\kern-.08em
%    T\kern-.1667em\lower.7ex\hbox{E}\kern-.125emX}}
%\providecommand\BibTeX{{\rm B\kern-.05em{\sc i\kern-.025em b}\kern-.08em
%    \TeX}}
\providecommand\BibTeX{{B\kern-.05em{\sc i\kern-.025em b}\kern-.08em
    \TeX}}
\providecommand\Matlab{\textsc{Matlab}}

\newcommand\webpage[1]{\href{http://cs.utexas.edu/~menie482/#1}{[link]}}
\newcommand\extwebpage[1]{\href{#1}{[link]}}

%%%%%%%%%%%%%%%%%%%%%%%% End Helper Commands %%%%%%%%%%%%%%%%%%%%%%%%%%%

%%%%%%%%%%%%%%%%%%%%%%%%% Begin CV Document %%%%%%%%%%%%%%%%%%%%%%%%%%%%

\begin{document}
\makeheading{Shun Zhang}

\section{Contact Information}

% NOTE: Mind where the & separators and \\ breaks are in the following
%       table.
%
% ALSO: \rcollength is the width of the right column of the table
%       (adjust it to your liking; default is 1.85in).
%
\newlength{\rcollength}\setlength{\rcollength}{2in}%
%
\begin{tabular}[t]{@{}p{\textwidth-\rcollength}p{\rcollength}}
\href{http://cs.utexas.edu/}%
     {Department of Computer Science} & \\
\href{http://www.utexas.edu/}{The University of Texas at Austin} \\
1 University Station A8000   & \includegraphics[width=110px]{phone.png}\\
Austin, TX  78712     & \email{jensen.zhang@utexas.edu}\\
\url{http://cs.utexas.edu/~menie482/} & \\
\end{tabular}


\section{Research Interests}

Reinforcement learning, robotics, multiagent systems, convex optimization, human
cognition.

\section{Education}

\href{http://www.utexas.edu}{\textbf{University of Texas at Austin}},
Austin, TX
\begin{outerlist}

\item[] Integrated B.S./M.S. Program,
        \href{http://cs.utexas.edu/}
             {Computer Science},
             Jan. 2012 - May. 2015 (Expected)
        \begin{innerlist}
        \item Major G.P.A. 3.8. Overall G.P.A. 3.55.
        \item \hyperref[sec:masterThesis]{Master Thesis}.
        \end{innerlist}
\end{outerlist}
\vspace{.1in}
\href{http://iao.nuaa.edu.cn/}{\textbf{Nanjing University of Aeronautics and
Astronautics}},
Nanjing, China
\begin{outerlist}
\item[] Undergraduate program,
             Computer Science and Technology,
             Sep. 2009 - Dec. 2011
        \begin{innerlist}
        \item G.P.A. 3.8/5.0.
        \item Transferred to University of Texas at Austin in Jan. 2012.
        \end{innerlist}

\end{outerlist}

\section{Research Experience}

\textbf{Modular Reinforcement Learning} \hfill {Fall 2014 --- Spring 2015}
\label{sec:masterThesis}
\begin{innerlist}
\item[] Department of Computer Science and Center for Perceptual Systems\\
        University of Texas at Austin
\vspace{.05in}
\item Supervisor: Prof. \href{http://www.cs.utexas.edu/~dana/}{Dana Ballard} and
Prof. \href{http://www.utexas.edu/cola/centers/cps/faculty/mmh739}{Mary Hayhoe}.
\item Research question: {\em Assume human already has Markov Decision Processes
(MDP) trained for preliminary tasks, how would these MDPs contribute to complicated
behaviors?}
\item Using Inverse Reinforcement Learning to interpret human's behavior,
assuming that it is a combination of the MDPs for preliminary tasks.
\end{innerlist}
\vspace{.1in}

\textbf{Determining Placements of Influencing Agents in a Flock} \hfill {Fall 2014}
\begin{innerlist}
\item[] Department of Computer Science\\
        University of Texas at Austin
\vspace{.05in}
\item Supervisor: Prof. \href{http://www.cs.utexas.edu/~pstone/}{Peter Stone}.
\item Research question: {\em Where should influencing agents be located within
a flock to maximize their influence on the flock?}
\item Using MASON simulator to evaluate different placements of influencing
agents, including border, grid, and graph-based placements.
\end{innerlist}
\vspace{.1in}

\textbf{Semi-Autonomous Intersection Management} \hfill {Summer, Fall 2012}
\begin{innerlist}
\item[] Department of Computer Science\\
        University of Texas at Austin
\vspace{.05in}
\item Supervisor: Prof. \href{http://www.cs.utexas.edu/~pstone/}{Peter Stone}
and Prof. \href{http://cse.unist.ac.kr/~chiu}{Tsz-Chiu Au}.
\item Research question: {\em Can we find a policy better than traffic signals,
if human-driven, semi-autonomous and fully-autonomous vehicles are sharing the
road?}
\item Designing and evaluating a policy that is competent with all three types
of vehicles, and performs better than traffic signals.
\end{innerlist}
\vspace{.1in}

\textbf{Reinforcement Learning on Atari Games} \hfill {Fall 2013}
\begin{innerlist}
\item Supervisor: Prof. \href{http://www.cs.utexas.edu/~pstone/}{Peter Stone}.
\item Undergraduate research course \webpage{papers/texplore.pdf}.
\end{innerlist}
\vspace{.1in}

\textbf{Action Selection in Robotic Motion Learning} \hfill {Fall 2013}
\begin{innerlist}
\item Supervisor: Prof. \href{http://www.cs.utexas.edu/~pstone/}{Peter Stone}.
\item Autonomous Robots course project. Achieved in Undergraduate Research
Journal in University of Texas at Austin, 2014
\extwebpage{http://texasurj.com/archives/2014.pdf}, page 77.
\end{innerlist}
\vspace{.1in}

\textbf{Structured Exploration for Relational Reinforcement Learning} \hfill
{Spring 2013}
\begin{innerlist}
\item Supervisor: Prof. \href{http://www.cs.utexas.edu/~pstone/}{Peter Stone}.
\item Reinforcement Learning course project
\webpage{papers/structuredExploration.pdf}.
\end{innerlist}
\vspace{.1in}

\section{Publications}
\begin{innerlist}
\item Tsz-Chiu Au, {\bf Shun Zhang}, and Peter Stone. Semi-Autonomous
Intersection Management (Extended Abstract).  Autonomous Agents and Multiagent
Systems (AAMAS), 2014. \webpage{papers/Au14SemiAIM.pdf}
\item Tsz-Chiu Au, {\bf Shun Zhang}, and Peter Stone. Autonomous Intersection
Management for Semi-Autonomous Vehicles. In Handbook of Transportation, 2015.
\webpage{papers/Routledge15-Au.pdf}
\footnote{This is the extended version of the Semi-Autonomous Intersection
Management paper.}
\item Katie Genter, {\bf Shun Zhang}, and Peter Stone. Determining Placements of
Influencing Agents in a Flock.  Autonomous Agents and Multiagent Systems
(AAMAS), 2015.
\item {\bf Shun Zhang}, Matthew Tong, Mary Hayhoe, Dana Ballard. Modular Inverse
Reinforcement Learning on Human Motion. Reinforcement Learning and Decision
Making (RLDM), 2015.
\end{innerlist}

\section{Presentations}
\begin{innerlist}
\item Intersection Management with Constraint-Based Reservation Systems. 
Autonomous Robots and Multirobot Systems (ARMS), 2014.
\end{innerlist}

\section{Conference Attendance}
\begin{innerlist}
\item Autonomous Agents and Multiagent Systems (AAMAS), Paris, 2014.
\item AAAI Conference, Austin, TX, 2015.
\end{innerlist}

\section{Courses and Projects}
Graduate Level
\begin{innerlist}
\item Large Scale Optimization (EE 381V)
\item Markov Chain and Mixing Time (M 394C)\\
Final Project: \textit{Mixing Time in Reinforcement Learning Convergence
Analysis .  \webpage{papers/mix.pdf}}
\item Machine Learning (CS 391L)\\
Project reports:
\begin{innerlist}[before=\itshape,font=\normalfont]
\item Principal Component Analysis. \webpage{391L/project1.pdf}
\item Independent Component Analysis. \webpage{391L/project2.pdf}
\item Approximate Inference in Bayesian Networks. \webpage{391L/project4.pdf}
\item Reinforcement Learning. \webpage{391L/project5.pdf}
\item Genetic Algorithm. \webpage{391L/project6.pdf}
\end{innerlist}
\item Autonomous Robots (CS 393R)
\item Randomized Algorithms (CS 388R)
\item Reinforcement Learning (CS 394R)\\
Project reports:
\begin{innerlist}[before=\itshape,font=\normalfont]
\item N-armed bandit Problem. \webpage{394proj1/report.html}
\item Eligibility Traces. \webpage{394proj2/report.html}
\item Bootstrapping with Function Approximation. \webpage{394proj3/report.html}
\item Transfer Learning in Gridworld. \webpage{394proj4/report.pdf}
\end{innerlist}
\item Introduction to Mathematical Logic (CS 388L)
\item Automated Logical Reasoning (CS 395T)
\end{innerlist}
\vspace{.1in}

Undergraduate Level
\begin{innerlist}
\item Artificial Intelligence (CS 343)
\item Principles of Computer Systems (CS 439)
\item Automata Theory (CS 341)
\item Information Retrieval (CS 371R)
\item Programming Languages (CS 345)\\
Final Project: \textit{List Interpreter. \webpage{345final/Lisp.pdf}}
\item etc.
\end{innerlist}

\section{Awards}
Student Awards --- University of Texas at Austin
\begin{innerlist}
\item Louis E. Rosier Memorial Endowment Scholarship. \hfill 2013-2014
\end{innerlist}
\vspace{.1in}

Student Awards --- Nanjing University of Aeronautics and Astronautics
\begin{innerlist}
\item Department Scholarships. \hfill 2009-2011
\end{innerlist}

\section{Teaching Experience}

\textbf{Undergraduate Teaching Assistant (Proctor)} \hfill {Fall 2013, Spring 2014}
\begin{innerlist}
\item[] CS 301K Foundations of Logical Thought\\
        with Dr. Jacob Schrum\\
        Department of Computer Science,\\
        University of Texas at Austin
\end{innerlist}

\section{Industry Experience}

\textbf{SDE Intern at Amazon} \hfill {Summer 2014}
\begin{innerlist}
\item[] Seattle, WA
\vspace{.05in}
\item Created an internal tool that supports WebRTC in Firefox.
\item Modifying an internal tool for visualizing customer data, using Scala on
Spark.
\end{innerlist}
\vspace{.1in}

\textbf{SDE Intern at Semantic Designs} \hfill {Summer 2013}
\begin{innerlist}
\item[] Austin, TX
\vspace{.05in}
\item Integrating a GUI viewer to Smart Differencer (TM) tools. Using JavaCC for
parsing and Swing for GUI.
\end{innerlist}

\section{Languages}

\begin{innerlist}
\item Natural languages: Mandarin Chinese (native), English (fluent), Japanese
(preliminary).
\item Programming languages: Proficient in programming in Python, Octave/Matlab,
Java, C/C++; Familiar with Lisp, Oracle SQL, \LaTeX, Web Development Languages
(HTML, JavaScript, PHP), Perl, Scala.
\end{innerlist}




\end{document}

%%%%%%%%%%%%%%%%%%%%%%%%%% End CV Document %%%%%%%%%%%%%%%%%%%%%%%%%%%%%

%----------------------------------------------------------------------%
% The following is copyright and licensing information for
% redistribution of this LaTeX source code; it also includes a liability
% statement. If this source code is not being redistributed to others,
% it may be omitted. It has no effect on the function of the above code.
%----------------------------------------------------------------------%
% Copyright (c) 2007, 2008, 2009, 2010, 2011 by Theodore P. Pavlic
%
% Unless otherwise expressly stated, this work is licensed under the
% Creative Commons Attribution-Noncommercial 3.0 United States License. To
% view a copy of this license, visit
% http://creativecommons.org/licenses/by-nc/3.0/us/ or send a letter to
% Creative Commons, 171 Second Street, Suite 300, San Francisco,
% California, 94105, USA.
%
% THE SOFTWARE IS PROVIDED "AS IS", WITHOUT WARRANTY OF ANY KIND, EXPRESS
% OR IMPLIED, INCLUDING BUT NOT LIMITED TO THE WARRANTIES OF
% MERCHANTABILITY, FITNESS FOR A PARTICULAR PURPOSE AND NONINFRINGEMENT.
% IN NO EVENT SHALL THE AUTHORS OR COPYRIGHT HOLDERS BE LIABLE FOR ANY
% CLAIM, DAMAGES OR OTHER LIABILITY, WHETHER IN AN ACTION OF CONTRACT,
% TORT OR OTHERWISE, ARISING FROM, OUT OF OR IN CONNECTION WITH THE
% SOFTWARE OR THE USE OR OTHER DEALINGS IN THE SOFTWARE.
%----------------------------------------------------------------------%
